\clearpage
\section{R2 robust-Sortino development strategy}
\label{sec:r2_update}

R2 is a separately versioned retrospective development strategy. It does not alter the
frozen R1 or R1.1 sources, artifacts, or freeze manifests, and it replaces R1.1 only if
every prespecified historical and simulation gate passes. Its Sortino target is zero net
return: because cash earns zero in this comparison, every negative post-cost month enters
downside deviation.

At each decision, R2 estimates two complete Greek factor/residual joint covariances. The
recent estimate uses the last 36 observed monthly returns; the expanding estimate uses all
prior observations from February 2015. Candidate recent weights of 0.25, 0.50, and 0.75
are selected by inner chronological covariance QLIKE. A one-standard-error rule selects
the lowest recent weight within the best estimate's uncertainty band, and the weight is
fixed at 0.50 until 12 inner forecasts exist. Both estimates retain the factor/residual
cross block; their joint covariance matrices are blended before transformation to option
space.

Missing option returns are not set to zero. A contract requires at least 24 observations
in the recent 36-date window. Missing cells receive the contemporaneous Greek-factor
prediction, and five deterministic scenario panels add residual draws resampled from that
contract's observed residuals. R2 evaluates the worst downside deviation across the recent
panel, an expanding six-month circular-block panel, and all five imputation panels. The
structural R1 premia channels remain in place, but the underlying and volatility premia use
expanding exponential weights with a 36-month half-life; the direct historical option-mean
weight remains zero.

Daily data condition risk, not option P\&L. For each surviving underlying, training data
are limited to the last 756 daily observations, with at least 500 required. HAR realized
variance and EWMA forecasts are combined over weights from zero to one in quarter steps by
rolling QLIKE, again with a one-standard-error preference for the more stable EWMA model.
The 21-day variance ratio is clipped to $[0.67,1.50]$. One diagonal transform scales delta
factor standard deviations by the volatility ratio and centered-quadratic factors by the
variance ratio; vega and residual dimensions are unchanged. This single joint transform
preserves positive semidefiniteness and the factor/residual cross terms. Missing daily
history produces an audited unscaled-covariance fallback.

The first optimization stage is a robust SOCP. It maximizes expected return net of
direction-specific predictable costs subject to unit worst-family zero-target downside
deviation. It uses one signed weight vector with \verb|abs|/positive-part expressions;
there are no overlapping positive and negative legs that can manufacture gross. Because
the ratio determines direction but not size, a separate scalar stage normalizes the
direction and maximizes expected net log growth. Scale zero is admissible and is selected
when no positive scale has positive expected log growth.

The scalar stage applies a 10 percent annualized downside-deviation ceiling and the R1.1
25 percent annualized volatility ceiling. It also retains the 10 percent monthly CVaR,
20 percent named stress loss, 75 percent short-margin, 100 percent collateral, assignment,
liquidity, Greek, beta, and concentration limits. The worst training three-month and
six-month compounded losses may not exceed 15 and 20 percent, respectively. The continuous
target is converted directly to whole contracts by truncating counts toward cash. If that
exact book breaches any constraint, the selected position is cash; no substitute portfolio
is introduced, and the rejected book's original diagnostics remain in the abstention
ledger.

\begin{table}[H]
\centering
\caption{R2 chronological retrospective development replay, February 2018--April 2026.
Sharpe is secondary; the promotion decision also uses path survival and the locked
simulation suites described in the text.}
\label{tab:r2_development}
\resizebox{\textwidth}{!}{\begin{tabular}{lrrrrrr}
\toprule
Universe & $N$ & Geom. return & Sortino & Max DD & Abstain & Verdict \\
\midrule
larger & 93 & 0.006 & 0.479 & -0.043 & 0 & development\_survived \\
larger+VIX & 93 & -0.012 & -0.265 & -0.127 & 0 & development\_survived \\
orig & 93 & 0.074 & 1.983 & -0.120 & 1 & development\_survived \\
orig+VIX & 93 & 0.023 & 0.397 & -0.170 & 1 & development\_survived \\
\bottomrule
\end{tabular}
}
\end{table}

The VIX-40 overlay remains excluded from headline R2 returns until complete executable
OPRA event quotes exist. It therefore cannot overwrite either R2's realized path or the
original E1 absorbed-zero finding. R2's 5,000 six-month circular-block paths, 2,000
60-month repriced paths under both joint-GARCH blocks and Gaussian-copula sensitivity, and
200 refitted histories per universe are locked diagnostics. CPCV remains a dependence and
fragility diagnostic rather than a tradable or confirmatory claim. A separate prospective
freeze requires 36 untouched future monthly observations before any confirmatory statement.
